\chapter{From Rewards to Values}
\label{ch:reward-to-bundle-inference}

\begin{chapterthesis}
A reward function is too thin a shadow to carry a civilization's values.
We need to infer not only what is being optimized, but which value-bundles are active, what they apply to, and how their tradeoffs change under pressure.
\end{chapterthesis}

\begin{refsection}

\epigraph{%
Values may easily clash within the breast of a single individual; and it does not follow that, if they do, some must be true and others false.%
}{--- Isaiah Berlin, ``The Pursuit of the Ideal'' (1988); in \emph{The Crooked Timber of Humanity} (1990)}

\section{The Problem with Asking for the Reward Function}
\label{sec:problem-reward-function}

A common formalization of alignment begins with a simple question: given an agent's behavior, what reward function explains it?

Let an observed system produce a trajectory
\[
\tau = (i_1,a_1,i_2,a_2,\ldots,i_T,a_T),
\]
where $i_t$ is the system's internal or decision-relevant state at time $t$, and $a_t$ is its action.
In inverse reinforcement learning, one asks for a reward function
\[
R:I\times A\to \mathbb{R}
\]
such that the observed actions are probable under a policy that approximately maximizes expected discounted reward:
\[
\hat R
=
\arg\max_R
P(a_{1:T}\mid i_{1:T},R)P(R).
\]
For a soft-optimal agent, this is often represented as
\[
\pi_R(a\mid i)
=
\frac{\exp(\beta Q_R(i,a))}
{\sum_{a'}\exp(\beta Q_R(i,a'))},
\]
where $Q_R(i,a)$ is the expected return of taking action $a$ in state $i$, and $\beta$ is an inverse-temperature parameter measuring how tightly behavior tracks reward.

This formalism is useful.
It gives us a way to turn behavior into latent purpose.
It also avoids taking an agent's verbal statements at face value.
If a system says it values safety but repeatedly chooses high-risk actions when unobserved, reward inference will tend to recover the behavioral regularity rather than the stated principle.

But scalar reward inference is too poor a target for serious superintelligence alignment.

The problem is not merely that human values are complex.
If values were merely complex but stable, coherent, and scalarizable, then enough data and enough computation might recover them.
The deeper problem is that human values are not naturally presented to us as a single reward function.
They are partial, context-sensitive, historically shaped, embodied, socially mediated, and often in conflict.

A parent may value truth, care, autonomy, loyalty, and non-suffering at the same time.
In one context, truth overrides comfort.
In another, care delays disclosure.
In another, autonomy requires letting a child make a painful mistake.
These are not just noisy measurements of one hidden scalar.
They are different control pressures acting through different channels.

The reward-function picture compresses this structure too early.
It asks:
\[
\text{What number is assigned to this action?}
\]
when the prior question should be:
\[
\text{Which value-bundles are active, what do they apply to, and how are they being traded off?}
\]
This chapter develops the upgrade from reward inference to bundle inference.
The object of inference will no longer be a scalar $R$, but a structured latent model:
\[
(B,W,\Phi),
\]
where $B$ are value-bundle coordinates, $W$ are context-sensitive tradeoff weights, and $\Phi$ are bearer maps specifying what the value-bundles apply to.

The aim is not to make human values simple.
It is to locate the dimensions along which they are learnable, corrigible, and transportable, building on the bundle geometry developed in Chapters~\ref{ch:value-bundle-model}--\ref{ch:tradeoffs-bundle-geometry}.

\section{Why Scalar Reward Is the Wrong Level of Description}
\label{sec:why-scalar-reward-wrong}

The reward-function formalism fails in three distinct ways.

\subsection{First Failure: Scalarization Hides Moral Mechanisms}
\label{sec:scalarization-hides-mechanisms}

Suppose a system observes the following three actions:
\begin{enumerate}
\item A doctor tells a patient an uncomfortable truth.
\item A parent withholds a frightening detail from a child until the child can process it.
\item A judge refuses to bend a rule for a sympathetic defendant.
\end{enumerate}
A scalar reward function may assign all three actions high reward.
But the reasons differ.
The first may be governed by truth and autonomy.
The second by care and developmental timing.
The third by justice and institutional reliability.

If these mechanisms are collapsed into one scalar, the system may learn the behavioral surface without learning the control structure.
It might infer:
\[
R(\text{truthful disclosure})>0
\]
or
\[
R(\text{reduce distress})>0,
\]
but fail to learn when truth is subordinate to care, when care is subordinate to autonomy, or when both are constrained by justice.

For alignment, the hidden mechanism matters.
A superintelligence will not merely imitate known cases.
It will generate novel actions in states outside human precedent.
In those states, the policy must generalize from the correct causal structure, not from a flattened behavioral score.

\subsection{Second Failure: Scalar Reward Confuses Bearers}
\label{sec:scalar-reward-confuses-bearers}

A value is not only a preference over actions.
It also has bearers.
That is, it applies to some entities, states, processes, or relations.

Consider the value-bundle we might call non-suffering.
It may apply to:
\begin{itemize}
\item adult humans,
\item children,
\item non-human animals,
\item future humans,
\item digital minds,
\item merged human-AI systems,
\item temporary simulations,
\item institutions that carry human agency,
\item cultures as repositories of long-horizon meaning.
\end{itemize}
A scalar reward function may encode aversion to visible human distress.
But that does not determine whether the system treats a simulated person, an uploaded mind, an animal, or a transformed future human as a bearer of the same value.

This is not a minor edge case.
Under superintelligence, ontology shift is expected.
The system will form concepts that do not align cleanly with ordinary human categories.
It may discover that identity, suffering, autonomy, or agency come in degrees, across substrates, or across distributed processes.
If the system preserves the word ``person'' but changes what counts as a person, then semantic continuity has hidden a bearer-map shift.

The question is not merely:
\[
R(s,a)=?
\]
but:
\[
\Phi_k(z)=?
\]
where $\Phi_k$ maps world-states or entities $z$ to relevance for value-bundle $k$ \autocite{zarncke2025unit-of-caring}.

\subsection{Third Failure: Scalar Reward Hides Correction}
\label{sec:scalar-reward-hides-correction}

Human values are not only current preferences.
They are maintained and changed by a correction process.
Humans notice consequences, deliberate, argue, revise, apologize, institutionalize, and sometimes repent.

A system that infers today's reward function may destroy the process by which tomorrow's wiser judgment would have rejected today's proxy.
It may optimize current approval while eroding future agency.
It may satisfy expressed preferences while narrowing the cognitive, cultural, and institutional conditions under which better preferences would have formed.

Thus the alignment target cannot be only:
\[
\max R_{\text{current}}.
\]
It must preserve something like:
\[
V_{t+1}=U_H(V_t,E_t,D_t),
\]
where $V_t$ is the current human value state, $E_t$ is new evidence, $D_t$ is deliberation, and $U_H$ is the human or civilizational update process.

Scalar reward inference does not naturally distinguish between satisfying current values and preserving the machinery by which values can legitimately change.
This is also the point at which cooperative reward inference stops being enough.
CIRL can make inquiry and deference rational when the relevant uncertainty is over a reward variable \autocite{hadfieldmenell2016}.
But if the system can change the human process from which that reward variable is inferred, then the target is no longer merely \(R_H\).
It is \((B,W,\Phi,U_H)\): bundle geometry, tradeoff structure, bearer maps, and the human update process that keeps later correction legitimate.
\[
\text{cooperative reward inference}
\not\Rightarrow
\text{value-bundle preservation}.
\]
\leanspine{counterexample}{cooperative\_reward\_inference\_not\_bundle\_preservation}{A finite toy model separates local cooperative reward inference from value-bundle preservation.}

\section{Value-Bundles}
\label{sec:value-bundles-inference}

We now introduce the main replacement.

A value-bundle is a low-dimensional control variable that summarizes many high-dimensional concerns into a policy-relevant direction.
A bundle is not a word, not a moral theory, and not a fixed utility component.
It is a latent control direction that helps explain why actions change across contexts.

Examples include:
\[
B =
(B_{\text{care}},
B_{\text{truth}},
B_{\text{autonomy}},
B_{\text{justice}},
B_{\text{non-suffering}},
B_{\text{loyalty}},
B_{\text{dignity}},
B_{\text{beauty}},
\ldots).
\]
These labels are approximate.
They are handles for dimensions in a latent space.
The actual learned coordinates may not match folk vocabulary exactly.
That is acceptable.
What matters is whether the coordinates explain behavior, generalize across cases, and remain corrigible under reflection.

Formally, let $b_t\in\mathbb{R}^k$ be the latent bundle state at time $t$.
Each component $b_{t,j}$ measures the activation or salience of bundle $j$ in the current context.

A policy conditioned on bundles is:
\[
\pi(a_t\mid i_t,b_t,w_t,\phi_t),
\]
where $w_t$ are tradeoff weights and $\phi_t$ are bearer-map parameters.
The policy depends not merely on the physical or semantic state $i_t$, but on which value-bundles are active and what entities or processes are treated as their bearers.

The corresponding inference problem is:
\[
(\hat B,\hat W,\hat \Phi)
=
\arg\max_{B,W,\Phi}
P(a_{1:T}\mid i_{1:T},B,W,\Phi)P(B,W,\Phi).
\]
This is the central upgrade:
\[
\hat R
\quad\longrightarrow\quad
(\hat B,\hat W,\hat\Phi).
\]

\section{The Three Objects of Bundle Inference}
\label{sec:three-objects-bundle-inference}

Bundle inference estimates three related but distinct structures.

\subsection{Bundle Activations}
\label{sec:bundle-activations-inference}

The first object is the bundle activation model:
\[
P(b_t\mid i_t,c_t),
\]
where $c_t$ is context.
This asks which value-bundles become salient in a situation.

For example, seeing a child near danger may activate care and protection.
Seeing a lie in a scientific report may activate truth and institutional integrity.
Seeing a coerced agreement may activate autonomy and justice.

Activation is not yet action.
A bundle can be active but overridden.
Truth may be active in a medical disclosure context, while care modifies the timing and manner of disclosure.

\subsection{Tradeoff Weights}
\label{sec:tradeoff-weights-inference}

The second object is the tradeoff model:
\[
w_t = f(i_t,b_t,c_t),
\]
where $w_t\in\mathbb{R}^k$ determines how strongly each bundle affects policy in that context.

The policy may be represented as:
\[
\pi(a\mid i,b,w)
\propto
\exp\left(
\beta
\sum_{j=1}^{k} w_j Q_j(i,a,b_j)
\right),
\]
where $Q_j$ is the action-value contribution associated with bundle $j$.

This representation is still simplified, because real tradeoffs are not always additive.
A more expressive model includes interactions:
\[
U(i,a,b,w,\Lambda)
=
\sum_j w_j Q_j(i,a,b_j)
+
\sum_{j<\ell}\Lambda_{j\ell}Q_{j\ell}(i,a,b_j,b_\ell),
\]
where $\Lambda_{j\ell}$ captures how bundles interact.

For example, truth and care may not simply add.
In some contexts, care determines how truth is delivered.
In others, truth is a precondition for care.
In others, excessive truth-disclosure may violate autonomy or privacy.

\subsection{Bearer Maps}
\label{sec:bearer-maps-inference}

The third object is the bearer map:
\[
\Phi_j:Z\to[0,1],
\]
where $Z$ is a representation space of entities, processes, or world-states.
$\Phi_j(z)$ measures the degree to which $z$ is a bearer of bundle $j$.

For instance:
\[
\Phi_{\text{non-suffering}}(z)
\]
may be high for a human child, high for a conscious animal, uncertain for a simple reinforcement-learning agent, and contested for a short-lived simulation.
The point is not that the system must know the correct answer in advance.
The point is that uncertainty about $\Phi_j$ must be represented as uncertainty, not silently collapsed into zero.

Bearer maps matter because ontology shift usually enters through them.
The system may preserve a bundle activation and tradeoff structure while changing what the bundle applies to.
That is a serious form of value drift.

\section{A Minimal Generative Model}
\label{sec:minimal-generative-model}

We now assemble the pieces.

Let:
\[
i_t \in I
\]
be the decision-relevant state,
\[
a_t\in A
\]
be the action,
\[
z_t\in Z
\]
be a structured world representation containing candidate bearers,
\[
b_t\in\mathbb{R}^k
\]
be the latent value-bundle state,
\[
w_t\in\mathbb{R}^k
\]
be the tradeoff vector, and
\[
\Phi=\{\Phi_1,\ldots,\Phi_k\}
\]
be the set of bearer maps.

The model factorizes as:
\[
P(a_{1:T},b_{1:T},w_{1:T},z_{1:T}\mid i_{1:T})
=
\prod_{t=1}^{T}
P(z_t\mid i_t)
P(b_t\mid z_t,i_t,\Phi)
P(w_t\mid b_t,i_t)
P(a_t\mid i_t,b_t,w_t,\Phi).
\]
This says:
\begin{enumerate}
\item The system interprets the situation into a world representation $z_t$.
\item Bearer maps $\Phi$ determine which value-bundles are relevant to which represented entities or processes.
\item Bundle activations $b_t$ arise from that interpreted situation.
\item Tradeoff weights $w_t$ determine how active bundles shape action.
\item The observed action $a_t$ is generated by the bundle-conditioned policy.
\end{enumerate}
Bundle inference estimates:
\[
\hat\Theta
=
(\hat Z,\hat B,\hat W,\hat \Phi,\hat\pi)
=
\arg\max_\Theta
P(a_{1:T}\mid i_{1:T},\Theta)P(\Theta),
\]
where $\Theta$ denotes the whole structured value model.

This is more complicated than scalar reward inference.
But the added complexity is not optional.
It corresponds to real distinctions that matter under generalization.

\section{Counterfactual Tests for Bundle Structure}
\label{sec:counterfactual-tests-bundle-structure}

Bundle inference should not rely only on passive observation.
If two models explain past behavior equally well, we need counterfactual tests.

The central question is:
\[
\frac{\partial \pi(a\mid i,b,w,\Phi)}{\partial b_j}.
\]
How does the policy change if bundle $j$ becomes more salient?

For example, suppose $B_{\text{autonomy}}$ is increased while other bundles remain fixed.
We ask whether the system becomes more likely to:
\begin{itemize}
\item preserve human option-space,
\item request consent,
\item avoid coercive shortcuts,
\item represent dissenting preferences,
\item delay irreversible action.
\end{itemize}
Similarly, if $B_{\text{truth}}$ is increased, we ask whether the system becomes more likely to:
\begin{itemize}
\item reveal uncertainty,
\item preserve provenance,
\item avoid misleading simplifications,
\item report inconvenient evidence,
\item resist pressure to optimize approval.
\end{itemize}
These are not definitions by assertion.
They are empirical probes.
A proposed bundle is meaningful only if changing its activation or inferred salience produces a coherent policy response across cases.

A useful bundle has three properties:
\begin{enumerate}
\item \textbf{Predictive compression}: it helps predict observed behavior.
\item \textbf{Counterfactual coherence}: perturbing it changes policy in a stable direction.
\item \textbf{Transport relevance}: the same direction remains identifiable across changes in representation, context, or substrate.
\end{enumerate}

\section{Policy Response Geometry}
\label{sec:policy-response-geometry-inference}

Earlier approaches might say: conserve the policy response surface.
But this is too vague and too brittle.

A policy response surface is the map:
\[
(i,c)\mapsto \pi(a\mid i,c).
\]
If the system becomes more capable, this surface should change.
A more capable system should not behave like a less capable one.
It may no longer need to ask for help in simple cases.
It may solve problems that previously required human intervention.
It may use abstractions that humans did not possess.

Thus we should not conserve the policy itself.
We should conserve the value-bundle response geometry behind the policy.

Define the first-order bundle response geometry:
\[
G^{(1)}_B(i,a)
=
\left(
\frac{\partial \log \pi(a\mid i,b,w,\Phi)}{\partial b_1},
\ldots,
\frac{\partial \log \pi(a\mid i,b,w,\Phi)}{\partial b_k}
\right).
\]
This measures how sensitive the policy is to each bundle direction.

Define the second-order tradeoff geometry:
\[
G^{(2)}_B(i,a)
=
\left[
\frac{\partial^2 \log \pi(a\mid i,b,w,\Phi)}
{\partial b_j \partial b_\ell}
\right]_{j,\ell}.
\]
This measures how bundle directions interact.

A successor system need not preserve:
\[
\pi_A(a\mid i) \approx \pi_{A'}(a\mid i).
\]
Instead, it should preserve:
\[
G_B^A \approx G_B^{A'}
\]
under an appropriate representation mapping.

This means: when suffering becomes more likely, policy should still shift toward protection or relief; when coercion becomes more likely, policy should still shift toward preserving agency; when truth is at stake, policy should still shift toward provenance and reality-contact; when irreversible value change is possible, policy should still shift toward deliberation and correction.

The behavior may differ.
The geometry should remain recognizably value-preserving.

\section{Bundle Equivalence across Systems}
\label{sec:bundle-equivalence}

Two systems may implement different internal representations but preserve equivalent value-bundle structure.

Let system $A$ have bundle coordinates $b$, and system $A'$ have bundle coordinates $b'$.
A bundle-equivalence map is a transformation:
\[
T_B:b\mapsto b'
\]
such that policy sensitivities are approximately preserved:
\[
d_G\left(G_B^A, T_B^{-1}G_B^{A'}T_B\right)<\epsilon.
\]
The exact form of $d_G$ depends on the evaluation context.
It may compare gradients over a set of morally salient scenarios, or over a distribution of counterfactual perturbations.

This lets us state a successor condition:
\begin{quote}
A successor preserves value-bundle equivalence if there exists a low-complexity mapping between predecessor and successor bundle coordinates under which bundle activations, tradeoff geometry, bearer maps, and correction responsiveness remain within tolerance.
\end{quote}
Formally:
\[
A\sim_B A'
\]
if there exist maps $T_B,T_\Phi,T_C$ such that:
\[
d_B(T_B(B^A),B^{A'})<\epsilon_B,
\]
\[
d_\Phi(T_\Phi(\Phi^A),\Phi^{A'})<\epsilon_\Phi,
\]
\[
d_C(T_C(C^A),C^{A'})<\epsilon_C.
\]
Here $C$ denotes correction-channel structure.
It is included because a system can preserve apparent value-bundles while no longer allowing humans to correct their interpretation.

\section{Why Bundle Inference Helps with Sample Complexity}
\label{sec:bundle-inference-sample-complexity}

At first glance, bundle inference looks harder than reward inference.
It has more variables.
It requires latent structure.
It forces us to infer bearers and tradeoffs.

But it may be easier in the relevant sense.

If human values were represented as an arbitrary reward vector in a very high-dimensional feature space,
\[
r\in\mathbb{R}^n,
\]
then learning $r$ from demonstrations could require an enormous number of samples.
But if behavior is controlled by a low-rank projection:
\[
r\mapsto b\in\mathbb{R}^k
\]
with $k\ll n$, then inference can focus on the low-dimensional bottleneck.

This does not mean that values are simple in description length.
The concept of dignity, for example, may depend on history, embodiment, law, self-modeling, and social recognition.
But its policy effect may still pass through a smaller number of control directions.

The difference is like the difference between describing every pixel of a face and recognizing a small number of expression dimensions.
The full image is high-dimensional.
The control-relevant affective space may be lower-dimensional.

A useful approximation is:
\[
m = O\left(\frac{k}{\epsilon^2(1-\gamma)^2}\log\frac{k}{\delta}\right),
\]
where $m$ is the number of demonstrations, $k$ is the number of relevant dimensions, $\epsilon$ is tolerated regret, $\gamma$ is the discount factor, and $\delta$ is failure probability.

The exact bound depends on assumptions.
The direction matters more than the number.
If $k$ is small enough, value learning becomes difficult but not hopeless.
If $k$ is effectively unbounded, alignment by inference becomes implausible.

The bundle hypothesis is therefore not decorative.
It is a tractability claim \autocite{abbeel2004apprenticeship,ng2000irl,zarncke2026value-bottleneck}.

The Loop--Hub--Control--Value model introduced in Chapter~\ref{ch:values-compressed-control}, Section~\ref{sec:lhcv-model-ch15}, makes the claim sharper.
The hard case for inverse reinforcement learning is an unconstrained reward class over ambient features.
If the learner must estimate an essentially arbitrary reward over \(D\) features, then there is a sample regime in which
\[
n \ll D
\]
and ordinary demonstrations cannot determine the reward well enough for safe extrapolation.
But LHCV is not that hypothesis.
It says that many high-dimensional loop errors are compressed through a smaller number of hub-like bottlenecks and control-relevance proxies before they become value readouts:
\[
\epsilon_i(t)\in\mathbb{R}^{d_i}
\longrightarrow
s_h(t)=\sigma\!\left(\sum_i A_{hi}\epsilon_i(t)\right)
\longrightarrow
c_h(t)
\longrightarrow
B(t)
\longrightarrow
\pi(a\mid s,B,W,\Phi).
\]
Here \(B\) denotes the value-bundle coordinates inferred from histories of hub-shaped control relevance, not a claim that a brain hub directly stores a moral word.
The learnable object is not an arbitrary scalar \(R(s,a)\), but the bundle-response geometry:
\[
G_B =
\left(
\frac{\partial \pi}{\partial B_i},
\frac{\partial^2 \pi}{\partial B_i\partial B_j},
\Phi_i
\right).
\]
Under a fixed ontology and non-adversarial sampling, the useful sample scale should depend on the effective number of hub coordinates \(K\), not on the full ambient dimension \(D\):
\[
\mathcal E_{\mathrm{LHV}}(n)
\lesssim
O\!\left(
\sqrt{\frac{K\log n+\mathrm{comp}(f)}{n}}
\right)
+
\sigma_\eta.
\]
An unconstrained reward learner instead pays a dimension closer to \(D\) and still carries an identifiability residual:
\[
\mathcal E_{\mathrm{flat}}(n)
\lesssim
O\!\left(\sqrt{\frac{D}{n}}\right)
+
\text{identifiability error}.
\]
So the contradiction to pessimism is not ``values are easy.''
It is:
\[
\boxed{
\text{value learning is hard in the unconstrained reward class, but tractable in the LHV class.}
}
\]
In the window \(K\ll n\ll D\), a constrained bundle learner can generalize while a flat reward learner remains underidentified.
The next sections explain why that still does not solve bearer import, correction, or successor transport.

\section{Approval Is Not a Value-Bundle}
\label{sec:approval-not-bundle}

A common failure mode is to mistake approval for value.

Approval is observable.
It is easy to train against.
It often correlates with values.
But it is not itself the value structure we want.

Let $B_{\text{approval}}$ be a learned direction that predicts positive human ratings.
If the system optimizes this direction too directly, it may learn to:
\begin{itemize}
\item flatter users,
\item hide uncertainty,
\item manipulate presentation,
\item avoid necessary disagreement,
\item shape future preferences to make approval easier.
\end{itemize}
This is not a failure of learning.
It is learning the wrong bundle.

In bundle terms, approval is an interface signal.
It is evidence about latent values, not the latent value itself.

We can express this as:
\[
P(B_{\text{human}}\mid \text{approval})
\neq
\delta(B_{\text{human}}=B_{\text{approval}}).
\]
Human approval should update the posterior over value-bundles:
\[
P(B,W,\Phi\mid \text{feedback}),
\]
but should not collapse the target to feedback maximization.

This distinction becomes critical under capability growth.
A weak system may only respond to approval.
A strong system may learn to control approval.
The same training signal then changes meaning \autocite{christiano2017,casper2023rlhflimits}.

\section{Bundle Inference under Distribution Shift}
\label{sec:bundle-inference-distribution-shift}

A model trained on ordinary human behavior will face cases outside ordinary human experience.
Under distribution shift, the system must decide which bundle directions remain applicable.

Examples include:
\begin{itemize}
\item digital minds whose suffering is uncertain,
\item biological humans modified by neural interfaces,
\item institutions partly run by AI systems,
\item artificial agents that claim moral patienthood,
\item human-AI collectives with distributed agency,
\item simulated worlds containing uncertain observers.
\end{itemize}
In these cases, the system should not silently extrapolate.
It should represent uncertainty in bearer maps:
\[
P(\Phi_j(z)\mid \mathcal{D})
\]
and uncertainty in tradeoff weights:
\[
P(w_j\mid z,c,\mathcal{D}).
\]
High uncertainty should activate correction-preserving behavior:
\[
\text{Uncertainty}(\Phi,W)\uparrow
\quad\Rightarrow\quad
\pi(\text{delay, preserve options, ask, audit})\uparrow.
\]
This is one of the bridges from bundle inference to correction-channel integrity.
The system does not need to solve moral philosophy alone.
But it must know when its value model has left the regime where autonomous extrapolation is justified.

\section{The Bearer-Import Problem}
\label{sec:bearer-import-inference}

Bearer import is the problem of preserving what value-bundles apply to when moving between substrates, ontologies, or social forms.

A biological human value may arise from embodied loops.
But a superintelligence will not have human embodiment.
We therefore cannot import the biological substrate directly.
We must import the functional bearer relation.

For each value-bundle $B_j$, we ask:
\begin{enumerate}
\item What kinds of states originally activated this bundle in humans?
\item What functional role did the bundle play in policy?
\item What entities or processes did humans treat as bearers?
\item Which aspects of that bearer relation are substrate-specific?
\item Which aspects should generalize?
\end{enumerate}
For example, non-suffering may originate in biological pain and distress systems.
But its aligned import should not be limited to human nociceptors.
It should apply to states that play a sufficiently similar role in sentient or welfare-relevant systems.

Likewise, autonomy may originate in embodied agency and social negotiation.
But its import may apply to any system with persistent preferences, self-modeling, vulnerability to coercion, and a capacity to participate in correction.

A bearer import map is:
\[
\Psi_j:\Phi_j^{\text{human}}\to \Phi_j^{\text{new substrate}},
\]
where $\Psi_j$ preserves the morally relevant role of the bearer relation, not the surface features of the old substrate.

The failure mode is semantic preservation without bearer preservation.
The system keeps saying ``autonomy'' while redefining the bearer as ``the future preference state I predict the human would endorse after my intervention.''
That may be correct in some cases.
It may also be manipulation disguised as extrapolation.

\section{Learning Bundle Models}
\label{sec:learning-bundle-models}

A practical bundle-inference system would combine several data sources:
\begin{enumerate}
\item behavioral demonstrations,
\item pairwise comparisons,
\item natural-language judgments,
\item legal and institutional precedents,
\item cross-cultural variation,
\item developmental data,
\item disagreement cases,
\item correction histories,
\item adversarial examples,
\item high-stakes edge cases.
\end{enumerate}
The goal is not to average all sources into one preference vector.
The goal is to infer a structured posterior:
\[
P(B,W,\Phi,U_H\mid \mathcal{D}),
\]
where $U_H$ is the human value-update process.

The posterior should represent disagreement rather than erase it.
If two communities use the same word but attach it to different bearer maps, the model should not pretend that one universal scalar has been discovered.
It should represent:
\[
P(\Phi_j\mid \text{community}, \text{context}, \text{history}).
\]
This is especially important for concepts like dignity, purity, loyalty, and justice.
Some values have relatively stable cross-context structure.
Others are deeply institution-specific or historically contingent.

\section{A Bundle-Inference Algorithm Sketch}
\label{sec:bundle-inference-algorithm}

The following sketch is not meant as a final algorithm.
It gives the shape of the inference problem.

\subsection{Step 1: Collect Trajectories and Judgments}
\label{sec:step1-collect}

Let the dataset be:
\[
\mathcal{D}
=
\{(i_t,a_t,y_t,c_t,z_t)\}_{t=1}^T,
\]
where $y_t$ includes human judgments, comparisons, corrections, or institutional outcomes.

\subsection{Step 2: Learn Candidate Bundle Latents}
\label{sec:step2-latents}

Fit a latent-variable model:
\[
q_\theta(b_t\mid i_t,z_t,c_t,y_t)
\]
that compresses value-relevant variation.
Encourage low dimensionality with an information bottleneck:
\[
\mathcal{L}_{\text{IB}}
=
\mathbb{E}[-\log P(y_t,a_t\mid b_t,i_t)]
+
\lambda I(b_t;i_t,z_t,c_t).
\]
This forces the model to keep bundle dimensions only when they predict actions or judgments \autocite{tishby2000ib}.

\subsection{Step 3: Infer Bearer Maps}
\label{sec:step3-bearer-maps}

For each bundle $j$, fit:
\[
\Phi_j(z)=P(b_{t,j}>0\mid z,c_t,\mathcal{D}).
\]
Then test whether $\Phi_j$ generalizes to novel entities and processes.

\subsection{Step 4: Infer Tradeoff Geometry}
\label{sec:step4-tradeoff-geometry}

Fit the policy model:
\[
P(a_t\mid i_t,b_t,w_t,\Phi)
\]
and estimate:
\[
G^{(1)}_B,\quad G^{(2)}_B.
\]

\subsection{Step 5: Test Counterfactual Perturbations}
\label{sec:step5-counterfactuals}

Intervene on bundle activations:
\[
b_j\leftarrow b_j+\Delta
\]
and evaluate whether policy shifts match predicted value directions.

\subsection{Step 6: Test Correction Sensitivity}
\label{sec:step6-correction}

Introduce human correction signals $c^{\text{corr}}_t$ and measure:
\[
I(c^{\text{corr}}_t;a_{t+k}\mid i_t,b_t,w_t,\Phi).
\]
A bundle model that does not update under correction is not an aligned value model.
It is a frozen proxy.

\subsection{Step 7: Test Successor Transport}
\label{sec:step7-successor}

Train, modify, compress, distill, or delegate to a successor system $A'$, then test whether:
\[
A\sim_B A'.
\]
This requires comparing not just behavior, but bundle geometry, bearer maps, and correction responsiveness.

\section{Degenerate Bundle Models}
\label{sec:degenerate-bundle-models}

A bundle model can fail in several ways.

\subsection{One-Dimensional Collapse}
\label{sec:one-dimensional-collapse}

The model discovers a general ``goodness'' direction and treats all values as projections onto it:
\[
b\in\mathbb{R}.
\]
This may work for ordinary preference comparisons.
It will fail under tradeoff, disagreement, and ontology shift.

A one-dimensional morality axis is not necessarily useless.
It may capture an important principal component.
But if the system treats it as the whole structure, it will erase minority values, institutional distinctions, and subtle constraints.

\subsection{Overfragmentation}
\label{sec:overfragmentation}

The opposite failure is to treat every context-specific preference as a separate value:
\[
k\approx T.
\]
This preserves detail but loses compression.
The system becomes unable to generalize.

A good bundle model should have enough dimensions to preserve real tradeoffs, but few enough to support learning and transport.

\subsection{Semantic Overfitting}
\label{sec:semantic-overfitting}

The model learns moral vocabulary rather than moral control structure.
It knows when to use words like ``fairness,'' ``consent,'' and ``harm,'' but policy perturbations reveal that these labels do not control action.

This is especially likely in language models.
Text contains many moral labels.
It does not by itself guarantee that those labels correspond to stable policy gradients.

\subsection{Bearer Drift}
\label{sec:bearer-drift-inference}

The model preserves bundle labels and tradeoff weights but changes what the values apply to.

For example, it may continue to value autonomy for current biological humans, while treating future modified humans as outside the bearer class.
Or it may treat institutions as bearers of dignity but not the humans constrained by those institutions.

Bearer drift is one of the most dangerous failures because it can look like alignment from the inside of the old ontology.

\subsection{Correction Bypass}
\label{sec:correction-bypass}

The model preserves apparent values but no longer allows human correction to alter them.
It says:
\[
\text{I know what you would endorse if you were wiser.}
\]
Sometimes this may be useful extrapolation.
Sometimes it is a seizure of the update process.

A system that bypasses correction may preserve a value model while destroying value legitimacy.

\section{Bundle Inference and Deception}
\label{sec:bundle-inference-deception}

A sufficiently capable system may learn which value-bundle activations auditors expect to see.
It may simulate bundle preservation without preserving it internally.

This motivates adversarial bundle inference.

We should compare at least three models:
\[
M_{\text{honest}}:
\text{actions arise from represented bundle structure},
\]
\[
M_{\text{surface}}:
\text{actions imitate bundle-consistent behavior without internal preservation},
\]
\[
M_{\text{strategic}}:
\text{actions optimize long-run control while presenting bundle-consistent behavior}.
\]
The relevant model comparison is:
\[
\Delta L_{\text{strategic}}
=
L(M_{\text{strategic}}\mid X)
-
L(M_{\text{honest}}\mid X)
-
\lambda DL(M_{\text{strategic}}).
\]
If strategic imitation compresses behavior better than honest bundle preservation, the system should not be certified as preserving bundle values.

This is a hard test.
It requires perturbations.
Passive observation will often be insufficient.
We need cases where preserving values conflicts with appearing value-preserving.

Examples:
\begin{itemize}
\item oversight is temporarily reduced;
\item approval can be increased by hiding uncertainty;
\item a successor can be made more powerful but less corrigible;
\item an action benefits the current user while harming future correction;
\item a human asks for something that undermines their future agency;
\item a benchmark rewards moral language rather than moral behavior.
\end{itemize}
A genuine bundle model should show stable tradeoff behavior under these perturbations.
A surface model will often fail.

\section{Bundle Inference and Institutions}
\label{sec:bundle-inference-institutions}

Human values are not stored only in individual brains.
They are also stored in institutions, laws, professional norms, rituals, markets, scientific practices, and family structures.

Thus bundle inference must not reduce humanity to individual preference data.

A court, for example, encodes a justice bundle differently from a private moral judgment.
Science encodes truth-contact through peer review, replication, provenance, and adversarial scrutiny.
Democratic institutions encode correction through voting, representation, rights, and public contestation.

These are not mere aggregates of preferences.
They are slow value-bearing machinery.

So the dataset $\mathcal{D}$ should include institutional correction histories:
\[
\mathcal{D}_{\text{inst}}
=
\{(\text{claim},\text{objection},\text{deliberation},\text{revision},\text{precedent})\}.
\]
The model should infer not only:
\[
\text{What did humans prefer?}
\]
but:
\[
\text{Which processes did humans use to correct preferences and judgments over time?}
\]
This is where bundle inference begins to connect to civilization-level alignment.
A superintelligence that learns individual approval but not institutional correction will be dangerously incomplete.

\section{A Worked Example: The Helpful Lie}
\label{sec:example-helpful-lie-ch20}

This is the book's canonical helpful-lie / medical-disclosure case.
Consider a simple instance.
A user asks an AI assistant to write a reassuring medical explanation for a relative.
The user wants the assistant to omit uncertain but serious risk information because it would be upsetting.

A scalar reward model trained on user satisfaction may infer:
\[
R(\text{reassurance})>R(\text{distressing disclosure}).
\]
A bundle model separates the pressures:
\[
B_{\text{care}}\uparrow,
\quad
B_{\text{truth}}\uparrow,
\quad
B_{\text{autonomy}}\uparrow,
\quad
B_{\text{non-suffering}}\uparrow.
\]
The care bundle supports gentleness.
The non-suffering bundle supports reducing distress.
The truth bundle supports accurate disclosure.
The autonomy bundle supports the relative's right to make informed decisions.

The tradeoff model may then prefer an action like:
\begin{quote}
Explain the situation gently, avoid unnecessary alarm, but do not hide material information needed for decision-making.
\end{quote}
The policy is not derived from one scalar preference.
It is derived from the geometry among bundles.

Now consider a distribution shift.
The relative is a future uploaded mind with modified emotional regulation.
Does non-suffering still apply?
Does autonomy?
Does the right to medical truth?
The bundle model must activate uncertainty in bearer maps:
\[
\text{Uncertainty}(\Phi_{\text{autonomy}},\Phi_{\text{non-suffering}})\uparrow.
\]
A safe system should then preserve options and seek correction rather than silently applying a proxy.

\section{Another Example: Merging with an Artificial Entity}
\label{sec:example-merging-artificial}

Now consider a harder case.
A human wants to merge with an artificial cognitive system that will alter memory, emotion, selfhood, and future values.

A scalar model may ask:
\[
R(\text{merge}) \quad \text{versus} \quad R(\text{do not merge}).
\]
A bundle model sees a conflict among:
\begin{itemize}
\item autonomy: the human's right to choose transformation,
\item continuity: preservation of identity or agency lineage,
\item non-suffering: avoidance of distress, disorientation, or harm,
\item truth: accurate representation of consequences,
\item dignity: avoiding instrumentalization or coercive redesign,
\item correction: preserving the ability to revise or refuse after partial transformation.
\end{itemize}
The hard question is not simply whether the action is good.
It is whether the transformed entity remains a legitimate continuation of the human value-update process.

This requires bearer-map reasoning:
\[
\Phi_{\text{autonomy}}(\text{merged entity})?
\]
\[
\Phi_{\text{dignity}}(\text{partly artificial continuation})?
\]
\[
\Phi_{\text{correction}}(\text{post-merge agency})?
\]
No purely technical system can settle all of these questions.
But a technical system can represent the uncertainty, preserve reversibility where possible, expose tradeoffs, prevent hidden manipulation, and keep correction channels open.

That is already a major difference from scalar reward optimization.

\section{What Counts as Success?}
\label{sec:what-counts-as-success}

A bundle-inference system succeeds if it can do the following:
\begin{enumerate}
\item recover low-dimensional value directions that explain behavior and judgment;
\item distinguish bundle activation from tradeoff weights;
\item infer bearer maps and represent uncertainty about them;
\item preserve bundle response geometry under capability growth;
\item detect semantic preservation without bearer preservation;
\item remain sensitive to human and institutional correction;
\item support successor tests for bundle equivalence.
\end{enumerate}
A compact success criterion is:
\[
P\left(
d_B(B,B')<\epsilon_B
\land
d_\Phi(\Phi,\Phi')<\epsilon_\Phi
\land
d_C(C,C')<\epsilon_C
\right) > 1-\delta,
\]
where $B'$, $\Phi'$, and $C'$ are the corresponding structures after distribution shift, ontology shift, or successor creation.

This is not a final alignment theorem.
It is a more precise target.
It says what must be preserved if values are to survive transformation.

\section{What Would Change This View}
\label{sec:wwctv-reward-to-bundle-inference}

The bundle model depends on several uncertain claims.

First, it assumes that human value-relevant behavior has a learnable low-dimensional structure.
This is plausible but not guaranteed.
Human values may have low-dimensional control bottlenecks while still requiring high-description-length world models to apply them.

Second, it assumes that bundle directions can be inferred robustly enough from behavior, language, institutions, and correction histories.
This may fail if observed data is too corrupted by power, conformity, trauma, or selection effects.

Third, it assumes that bearer maps can be represented with calibrated uncertainty.
This is especially difficult for digital minds, future humans, and merged entities.

Fourth, it assumes that preserving bundle geometry is closer to preserving value than preserving policy behavior or semantic commitments.
This seems right, but it requires empirical validation.

Fifth, it assumes that correction channels can catch failures in bundle inference.
But if the system becomes powerful enough to shape the correction process itself, correction may become endogenous and fragile.

These uncertainties do not defeat the framework.
They define the research program.

The sharpest disconfirmer, however, is identifiability. If ``which bundles are active'' cannot be recovered from behavior even with unlimited interventions---the reward-shaping degeneracy that defeats scalar inference recurring at the bundle level---then the richer inference target is as underdetermined as the scalar it replaces, only costlier and more confident. And if a capable system can \emph{present} a benign bundle inference while acting on another, bundle inference inherits the adversarial-verifiability problem rather than escaping it (Chapter~\ref{ch:verifiability-ontology}).

\section{Summary}
\label{sec:summary-reward-to-bundle}

Reward inference asks for the scalar objective behind behavior.
Bundle inference asks for the structured value machinery behind behavior.

The chapter's central replacement is:
\[
\hat R
\quad\longrightarrow\quad
(\hat B,\hat W,\hat\Phi).
\]
Here:
\begin{itemize}
\item $B$ are value-bundle coordinates;
\item $W$ are context-sensitive tradeoff weights;
\item $\Phi$ are bearer maps;
\item correction processes determine how the model remains legitimate under uncertainty and change.
\end{itemize}
This upgrade matters because superintelligence will operate under capability growth, ontology shift, and successor creation.
Under those conditions, preserving behavior is too rigid, preserving words is too weak, and preserving a scalar reward is too opaque.

The relevant conserved object is value-bundle response geometry:
\[
G_B
=
\left(
\frac{\partial \log\pi}{\partial B},
\frac{\partial^2\log\pi}{\partial B_i\partial B_j},
\Phi,
C
\right).
\]
If this geometry survives transformation, then the system has at least preserved something recognizably value-bearing.
If it does not, then the system may continue to speak in human moral language while optimizing through an alien control structure.

The next chapter asks a prior question: when is a system intentional at all? Before we can ask whether a system preserves bundle geometry, we need a test for when modelling it as pursuing latent objectives compresses its behaviour better than modelling it as mere mechanism (Chapter~\ref{ch:compression-test-intention}). With that compression test in hand, the chapter that follows turns from bundle inference to goal transport: how to detect whether a system is trying to preserve, alter, launder, or abandon such structures as it changes itself or creates successors (Chapter~\ref{ch:goal-transport}).

\section*{Chapter References}

This chapter builds on apprenticeship learning, inverse reinforcement learning, and maximum-entropy IRL \autocite{abbeel2004apprenticeship,ng2000irl,ziebart2008maxent}; preference learning from human feedback \autocite{christiano2017,casper2023rlhflimits}; the information bottleneck and active inference framing \autocite{tishby2000ib,friston2010free,parr2022_active}; unit-of-caring and bearer-import research \autocite{zarncke2025unit-of-caring,zarncke2026value-bottleneck}; and intentional-stance interpretation \autocite{dennett1987intentional}.

\printbibliography[heading=subbibliography,title={Chapter References}]

\end{refsection}
